% v2.7.0 figure UID: FIG-P580-01
% Analytic contract on X=[0,5]: p=6x(5-x)/125, q_L=(2/5)1_[0,5/2], q_R=1/5.
\pgfkeysifdefined{/tikz/alt}{}{\tikzset{alt/.code={}}}
\pgfmathdeclarefunction{slpivtarget}{1}{%
  \pgfmathparse{6*(#1)*(5-(#1))/125}%
}
\def\ISXMin{0}
\def\ISXMax{5}
\def\ISXCut{2.5}
\pgfmathsetmacro{\ISQLHeight}{2/5}
\pgfmathsetmacro{\ISQRHeight}{1/5}
\pgfmathsetmacro{\ISPAtOne}{slpivtarget(1)}
\pgfmathsetmacro{\ISPAtMid}{slpivtarget(2.5)}
\pgfmathsetmacro{\ISPAtFour}{slpivtarget(4)}
\pgfmathsetmacro{\ISWAtOne}{\ISPAtOne/\ISQRHeight}
\pgfmathsetmacro{\ISWAtMid}{\ISPAtMid/\ISQRHeight}
\pgfmathsetmacro{\ISWAtFour}{\ISPAtFour/\ISQRHeight}
\tikzset{slfig-FIG-P580-01/.style={font=\fontsize{9.6pt}{11.6pt}\selectfont,
  every path/.append style={line cap=round,line join=round}}}
\pgfplotsset{slfig-FIG-P580-01-axis/.style={
  axis lines=left,xmin=0,xmax=5,ymin=0,ymax=.61,
  xtick={0,1,2.5,4,5},xticklabels={$0$,$1$,$5/2$,$4$,$5$},
  ytick={0,.2,.3,.4,.5},yticklabels={$0$,$1/5$,$3/10$,$2/5$,$1/2$},
  tick label style={font=\fontsize{9.6pt}{11.6pt}\selectfont},
  label style={font=\fontsize{9.6pt}{11.6pt}\selectfont},
  title style={font=\fontsize{10.2pt}{12.2pt}\selectfont,yshift=1pt},
  axis line style={draw=SLGray,line width=.72pt},
  x axis line style={shorten >=-2mm},
  tick style={draw=SLGray,line width=.60pt},
  xlabel style={align=center},ylabel style={align=center},clip=false}}
\begin{figure}[htbp]
\centering
\begin{tikzpicture}[slfig-FIG-P580-01,
  alt={对象：共同定义域[0,5]上的目标密度p(x)=6x(5-x)/125、左提议密度qL(x)等于2/5乘以区间[0,5/2]的示性函数、右提议密度qR(x)=1/5；关系：左图在x=5/2处标出qL的支撑边界，并以纹理显示其后qL=0而p>0的缺失区；右图qR处处为正，曲线上的圆、方、三角点对应x=1、5/2、4，独立比率卡由同一密度公式复算三个真实比率；结论：重要性抽样要求p绝对连续于q，支持缺口不能由增加样本或有限加权恢复，而支持覆盖本身不推出低方差或估计可靠。}]
\begin{groupplot}[group style={group size=2 by 1,horizontal sep=10mm},
  slfig-FIG-P580-01-axis,width=6.45cm,height=5.9cm]
\nextgroupplot[title={支持不足\quad $p\not\ll q_L$},
  ylabel={密度\\[-.2mm]每 $x$ 单位},
  xlabel={共同定义域 $x$ 从 $0$ 到 $5$\\[.3mm]
    斜线区\enspace $\textstyle q_L(x)$ 为 $0$ 且
    $\textstyle p(x)$ 为正}]
  \addplot[name path=pL,draw=none,domain=\ISXMin:\ISXMax,samples=301]
    {slpivtarget(x)};
  \addplot[name path=zeroL,draw=none,domain=\ISXMin:\ISXMax,samples=2]{0};
  \addplot[pattern=north east lines,pattern color=SLTextGray,draw=none]
    fill between[of=pL and zeroL,soft clip={domain=2.5:5}];
  \addplot[draw=SLBlue,line width=1.15pt,domain=\ISXMin:\ISXMax,samples=301]
    {slpivtarget(x)};
  \addplot[draw=SLTeal,line width=.95pt,densely dashed,
    domain=\ISXMin:\ISXCut,samples=2]{\ISQLHeight};
  \addplot[draw=SLTeal,line width=.95pt,densely dashed,
    domain=\ISXCut:\ISXMax,samples=2]{0};
  \addplot[only marks,mark=square*,mark size=2.8pt,draw=SLTeal,fill=SLTeal]
    coordinates {(\ISXCut,\ISQLHeight)};
  \addplot[only marks,mark=o,mark size=2.8pt,draw=SLTeal,line width=.90pt]
    coordinates {(\ISXCut,0)};
  \draw[draw=SLGray,line width=.80pt,densely dotted]
    (axis cs:\ISXCut,0)--(axis cs:\ISXCut,.455);
  \node[anchor=south west,align=left,text=SLTeal,inner sep=0pt,
    font=\fontsize{9.6pt}{11.6pt}\selectfont]
    at (axis cs:.16,.455) {虚线 $\textstyle q_L(x)$\\[-.1mm]取 $2/5$};
  \node[anchor=south west,text=SLBlue,inner sep=0pt,
    font=\fontsize{9.6pt}{11.6pt}\selectfont]
    at (axis cs:.16,.305) {实线 $p(x)$};
  \node[anchor=south,align=center,text=SLTextGray,inner sep=0pt,
    font=\fontsize{9.6pt}{11.6pt}\selectfont]
    at (axis cs:3.82,.455) {点线\enspace 支撑边界\\[-.1mm]横坐标取 $5/2$};

\nextgroupplot[title={支持覆盖\quad $q_R$ 覆盖目标},
  xlabel={共同定义域 $x$ 从 $0$ 到 $5$\\[.3mm]
    \textcolor{SLBlue}{实线 $\textstyle p(x)$}\quad
    \textcolor{SLTeal}{虚线 $\textstyle q_R(x)$ 为 $1/5$}}]
  \addplot[draw=SLBlue,line width=1.15pt,domain=\ISXMin:\ISXMax,samples=301]
    {slpivtarget(x)};
  \addplot[draw=SLTeal,line width=.95pt,
    dash pattern=on 3pt off 10pt on 3pt off 2pt on 3pt off 2pt
      on 3pt off 2pt on 3pt off 2pt on 3pt off 2pt
      on 3pt off 2pt on 3pt off 2pt on 3pt off 2pt
      on 3pt off 2pt on 3pt off 2pt on 3pt off 2pt
      on 3pt off 2pt on 3pt off 2pt on 3pt off 2.1pt,
    dash phase=63.4pt,
    domain=\ISXMin:\ISXMax,samples=2]{\ISQRHeight};
  \addplot[only marks,mark=o,mark size=3.0pt,draw=SLBlue,line width=.90pt]
    coordinates {(1,\ISPAtOne)};
  \addplot[only marks,mark=square*,mark size=2.8pt,draw=SLBlue,fill=white,
    line width=.90pt] coordinates {(\ISXCut,\ISPAtMid)};
  \addplot[only marks,mark=triangle*,mark size=3.2pt,draw=SLBlue,fill=white,
    line width=.90pt] coordinates {(4,\ISPAtFour)};
  \path[draw=SLGray,rounded corners=1.5pt,fill=white,line width=.55pt]
    (axis cs:.18,.340) rectangle (axis cs:4.82,.590);
  \node[anchor=center,align=center,inner sep=0pt]
    at (axis cs:2.5,.465) {$
      \begin{gathered}
        w(x)\text{ 由 }p(x)/q_R(x)\text{ 计算}\\[-.1mm]
        \begin{array}{ccc}
          w(1) & w(5/2) & w(4)\\[-.1mm]
          24/25 & 3/2 & 24/25
        \end{array}
      \end{gathered}$};
\end{groupplot}
\end{tikzpicture}
\caption{重要性抽样要求 $p\ll q$\quad 其中 $q$ 未覆盖的目标区域不能由增加样本或有限加权恢复}\label{fig:V5-C02-is-support}
\end{figure}
